\section{Multi-Agent Reinforcement Learning}

% ---------------------------------------------------------------- Thread transition
\begin{frame}{}
\vfill
\centering
{\large Thread 2 of 3}\\[0.5em]
{\LARGE \textbf{Multi-Agent RL}}\\[0.8em]
{\footnotesize\itshape The RLHF deep dive is done. A shorter overview now: what
changes when \emph{other learning agents} share your environment.}
\vfill
\end{frame}

% ================================================================ 4.1
% Why MARL is hard (non-stationarity) - TikZ
\begin{frame}{Why Multi-Agent RL Is Hard: Non-Stationarity}
\centering
\resizebox{0.92\linewidth}{!}{%
\begin{tikzpicture}[font=\footnotesize,
    box/.style={draw, rounded corners, minimum width=2.2cm, minimum height=1.0cm, align=center},
    ar/.style={-{Stealth[length=2.2mm]}, thick}]

  % ---- single-agent ----
  \node[font=\bfseries] at (1.6,3.2) {Single-agent};
  \node[box, fill=sysblue] (a1) at (0,1.6) {Agent};
  \node[box, fill=boxgreen](e1) at (3.2,1.6){Environment\\\scriptsize stationary};
  \draw[ar] (a1.north) to[bend left=35] node[above,font=\scriptsize]{action} (e1.north);
  \draw[ar] (e1.south) to[bend left=35] node[below,font=\scriptsize]{state, reward} (a1.south);

  % divider
  \draw[gray, dashed] (5.4,-0.6) -- (5.4,3.4);

  % ---- multi-agent ----
  \node[font=\bfseries] at (9.5,3.2) {Multi-agent};
  \node[box, fill=sysblue] (a2) at (7.0,1.6) {Agent $i$};
  \node[draw, rounded corners, fill=boxred, align=center, minimum width=3.6cm, minimum height=1.7cm]
        (e2) at (11.0,1.6) {Environment\\\scriptsize \emph{contains other}\\\scriptsize \emph{learning agents}\\\scriptsize non-stationary};
  \draw[ar] (a2.north) to[bend left=25] node[above,font=\scriptsize]{action} (e2.north west);
  \draw[ar] (e2.south west) to[bend left=25] node[below,font=\scriptsize]{state, reward} (a2.south);
\end{tikzpicture}%
}
\vfill
\begin{itemize}\itemsep0.25em\small
    \item Each agent's environment is \textbf{non-stationary}: the others are
          updating their policies at the same time, so the ground shifts as they
          learn; single-agent convergence guarantees break.
    \item Formalized as a \textbf{Markov game} (stochastic game), the multi-agent
          generalization of the MDP.
    \item Three regimes along one axis: \textbf{cooperative} / \textbf{competitive}
          / \textbf{mixed}.
\end{itemize}
\end{frame}

